\section{MPI Measurements}
\label{sec:mpi-results}

The MPI timings used Nicholas Johnson's branch \texttt{feature/mpi-block-plu} at commit \texttt{12b7cd3}, cloned directly because that branch was not yet merged into the group \texttt{main}.
The field dependency on that branch is the same cached fork used in Section~\ref{sec:results}.
Jobs were submitted with \texttt{sbatch} to the course cluster's Slurm partition \texttt{demo}.
That partition allows 5 minutes and has four machines, with two CPUs on each.
The MPI library is OpenMPI, loaded on the compute nodes.
The block size was 32 for every run.
Sequential mode launches one rank and uses the one-rank elimination.
The eight-rank mode places two ranks on each of the four machines.
The driver still forms \(C = AB\) on rank 0 before the determinants, so the product time is almost the same in both modes.
A run is accepted when \(\det(C) = \det(A)\det(B)\).
That check passed for every size in Table~\ref{tab:mpi}.
At \(n = 512\) the three determinants printed by the sequential run and by the eight-rank run were the same field elements.

\begin{table}[!t]
\caption{Cached sequential time against eight MPI ranks on four machines. Determinant time is the sum of \(\det(A)\), \(\det(B)\), and \(\det(C)\).}
\label{tab:mpi}
\centering
\scriptsize
\begin{tabular}{@{}rrrrr@{}}
\toprule
\(n\) & Seq. (s) & 8 ranks (s) & Seq. dets (s) & 8-rank dets (s) \\
\midrule
64  & 1.76 & 2.61 & 0.25 & 0.21 \\
128 & 5.53 & 5.61 & 2.04 & 1.09 \\
256 & 35.0 & 27.4 & 15.6 & 5.66 \\
512 & 272  & 200  & 120  & 33.5 \\
\bottomrule
\end{tabular}
\end{table}

Eight ranks are already slightly faster on the determinants at \(n = 64\), 0.21\,s against 0.25\,s.
The whole program is still slower there, 2.61\,s against 1.76\,s, because the table fill and the sequential product are larger than the determinant savings.
At \(n = 128\) the determinants drop from 2.04\,s to 1.09\,s, while the whole program is essentially tied, 5.61\,s against 5.53\,s.
The smallest tested size at which the whole program is faster with eight ranks is \(n = 256\), 27.4\,s against 35.0\,s.
At \(n = 512\) the determinant portion falls from 120\,s to 33.5\,s, a factor of about 3.6 on eight ranks, and the whole program falls from 272\,s to 200\,s.
The product at that size is still about 151\,s sequential and 164\,s with eight ranks, so most of the remaining time is work this MPI code does not divide.
Figure~\ref{fig:mpi} plots the determinant column of Table~\ref{tab:mpi}.

\begin{figure}[!t]
\centering
\begin{tikzpicture}
\begin{axis}[
  ieee lines,
  ymode=log,
  ymin=0.1,
  ylabel={Determinant time (s)},
  xlabel={Matrix dimension \(n\)},
  symbolic x coords={64,128,256,512},
]
\addplot coordinates {(64,0.25) (128,2.04) (256,15.58) (512,119.70)};
\addplot coordinates {(64,0.21) (128,1.09) (256,5.66) (512,33.53)};
\legend{Sequential, 8 ranks}
\end{axis}
\end{tikzpicture}
\caption{Sum of the three determinant calls on the four-machine cluster. The vertical axis is logarithmic.}
\label{fig:mpi}
\end{figure}

A second sweep used two ranks on one machine, at \(n \in \{64, 128, 256\}\), with the same block size and the same cached field.
The determinant sums were 0.23\,s against 0.23\,s at \(n = 64\), 1.91\,s against 1.61\,s at \(n = 128\), and 15.9\,s against 12.6\,s at \(n = 256\).
The whole program was faster with two ranks only at \(n = 256\), 33.9\,s against 35.6\,s.
Two ranks therefore reach a whole-program win at the same tested size as eight ranks, but the determinant speedup is smaller.
The identity check passed on that sweep as well.

The cluster could not host a larger size.
The \(512 \times 512\) sequential run took 272\,s of a 300\,s job.
An attempt to follow it with the eight-rank run in the same job was cancelled at the time limit.
The eight-rank \(512 \times 512\) numbers in Table~\ref{tab:mpi} come from a separate job that ran only that mode.
The step from 35\,s at \(n = 256\) to 272\,s at \(n = 512\) is a factor of 7.8, against the factor of 8 predicted by \(n^3\).
Applying another factor of \((4000/512)^3 \approx 477\) to the 272\,s measurement puts a cached sequential \(4000 \times 4000\) run near 36 hours on these cluster CPUs.
That case was not run.
The assignment's 20-minute cached figure remains the Apple M3 Max note, measured on a different implementation and a different computer.
